\section{Euclid除法}

\begin{theorem}{带余除法的存在性和唯一性}{}
	设$a,b$是自然数,$b>0$,则存在唯一的自然数$q$和$r$使得$a=bq+r\wedge r<b$.
\end{theorem}

\begin{definition}{余数,商,整数部分}{}
	设$a,b,q,r$是自然数,$0<r<b,a=bq+r$.
	\begin{enumerate}
		\item 我们称$r$是$a$除以$b$的余数.当$r=0$时称$a$是$b$的倍数(或$a$能被$b$整除,$b$是$a$的因数/因子).
		\item 若$a$是$b$的倍数,则把$q$记作$\frac{a}{b}$(或$a/b$),称之为$a$除以$b$的商;若$a$不是$b$的倍数,则称$q$是$a$除以$b$的商的整数部分.
	\end{enumerate}
\end{definition}

\begin{proposition}{}{}
	设$a,b,c,d$是自然数.
	\begin{enumerate}
		\item 若$a$是$b$的倍数,$c$是$a$的倍数,则$c$是$b$的倍数.进一步,当$a\neq 0$时$\frac{c}{b}=\frac{c}{a}\cdot\frac{a}{b}$.
		\item 若$c$和$d$都是$b$的倍数,则$c+d$是$b$的倍数,并且$\frac{c+d}{b}=\frac{c}{b}+\frac{d}{b}$.进一步,当$d\leq c$时$c-d$是$b$的倍数,并且$\frac{c-d}{b}=\frac{c}{b}-\frac{d}{b}$.
	\end{enumerate}
\end{proposition}

\begin{definition}{奇数,偶数}{}
	能被$2$整除的自然数称为奇数,不能被$2$整除的自然数被称为偶数.
\end{definition}

\begin{corollary}{}{}
	每个奇数都能表示成$2n+1$,其中$n$是自然数.
\end{corollary}










































